%!TEX root = ../main.tex
\chapter{Liquid Neural Networks}
\label{ch:lnn}

\section{Overview}
\label{sec:lnn-overview}

The network this work trains takes the place of the pilot: it receives the
observation defined in Section~\ref{sec:observation-space} --- the local voxel
map, the pose, the goal and the declaration the robot makes about itself --- and
produces the velocity command the platform executes. The problem has three
properties that determine the architecture.

It is \emph{partially observable}: a single frame does not say whether the robot
was accelerating, whether the obstacle ahead is approaching, or how long it has
been standing still. Memory is therefore expected to help --- an expectation,
not a premise, which the stateless control of
Section~\ref{sec:networks-implementation} exists to test.

It is \emph{non-uniform in time}: the rate at which frames arrive is not
constant. A discrete-time recurrent network, such as an LSTM
\cite{hochreiter1997lstm} or a GRU \cite{cho2014gru}, treats every step as
equivalent to the previous one, and therefore interprets a two-second
interruption as though it were an interval of a tenth of a second. This is the
reason why the family of cells adopted is that of \emph{Liquid Neural Networks}
\cite{hasani2021ltc,hasani2022cfc}: the state evolves not by iteration but by
integration of a differential equation, and the elapsed time is an input of the
update rather than an implicit assumption. The same family has been employed
successfully in end-to-end learnt flight navigation \cite{chahine2023flight},
which shares with the present problem the nature of the input and the need to
generalise outside the training distribution.

It is, finally, \emph{heterogeneous}: the same network must fly different
platforms, with velocity limits that vary and masses that differ. The
architecture therefore contains two modules --- the observation normaliser and
the robot-model embedding --- whose task is to make comparable robots that are
not.

\paragraph{The structure.}
The architecture is a single one, with two points of variation
(Figure~\ref{fig:architecture}). Three
independent branches --- the sparse encoder that reduces the voxel map to a
compact descriptor, the normalisation of the \num{41} dense values, and the
robot-model embedding --- are concatenated and projected by a single linear
layer into a fusion vector. The fusion vector traverses two recurrent cells in
cascade, and the state of the second feeds a linear head that produces the six
components of the command, rescaled by a final affine stage. The two points of
variation are \emph{which cell} occupies the recurrent core --- LTC, CfC or
LRCU --- and \emph{whether that core is sparsely wired} according to the NCP
scheme. Everything else lives in a common base class and is identical across all
variants, so that a change to the input or output path reaches every variant at
once and the comparison isolates exactly the two choices.

\begin{figure}[htbp]
  \centering
  \linespread{1}\selectfont
  \begin{tikzpicture}[node distance=6mm]
    \node[blk] (enc) {sparse encoder\\\footnotesize 4 levels $\to$ 200};
    \node[blk, right=9mm of enc] (nrm) {observation\\normaliser\\\footnotesize 41};
    \node[blk, right=9mm of nrm] (emb) {model\\embedding\\\footnotesize id $\to$ 8};

    \node[lbl, above=2mm of enc] {voxel indices};
    \node[lbl, above=2mm of nrm] {goal $\mid$ odom $\mid$ state};
    \node[lbl, above=2mm of emb] {\texttt{model\_id}};

    \node[blk, below=9mm of nrm] (fus) {fusion layer\\\footnotesize $249 \to 192$};
    \draw[flow] (enc.south) |- ([yshift=4mm]fus.north) -- (fus.north);
    \draw[flow] (nrm) -- (fus);
    \draw[flow] (emb.south) |- ([yshift=4mm]fus.north) -- (fus.north);

    \node[blk, below=10mm of fus] (c1) {cell 1\\\footnotesize $192 \to 80$};
    \node[blk, below=9mm of c1]  (c2) {cell 2\\\footnotesize $80 \to 80$};
    \draw[flow] (fus) -- (c1);
    \draw[flow] (c1) -- (c2);

    \begin{scope}[on background layer]
      \node[draw, dashed, rounded corners=3pt, fill=black!3,
            fit=(c1)(c2), inner sep=4mm] (core) {};
    \end{scope}
    \node[note, left=8mm of core, text width=27mm, anchor=east]
      {the variable part:\\LTC, CfC or LRCU,\\dense or NCP-wired};

    \node[blk, below=10mm of c2] (head) {head\\\footnotesize $80 \to 128 \to 6$};
    \node[blk, below=8mm of head] (cal) {output calibration};
    \node[lbl, below=6mm of cal] (out) {\texttt{cmd\_vel}, 6};
    \draw[flow] (c2) -- (head);
    \draw[flow] (head) -- (cal);
    \draw[flow] (cal) -- (out);

    \draw[thin flow, dashed] (fus.east) -- ++(14mm,0) |- (c2.east);
    \draw[thin flow, dashed] (c1.east)  -- ++(7mm,0)  |- (head.east);
  \end{tikzpicture}
  \caption{The architecture, with the two points of variation. Everything
           outside the dashed box is common to all the networks compared; inside
           it sit the two recurrent cells, which differ in their update rule and
           in whether their matrices are masked. The dashed arrows on the right
           are the two bypass paths that only the sparsely wired variants carry
           (Section~\ref{sec:lnn-wiring}).}
  \label{fig:architecture}
\end{figure}

\section{Continuous-Time Cells: LTC, CfC and LRCU}
\label{sec:lnn-taxonomy}

\subsection{From neural differential equations to liquid cells}
\label{sec:lnn-background}

The three cells employed in this work belong to a single line of models, and
understanding them requires that line to be briefly retraced.

The point of departure is the family of \emph{neural ordinary differential
equations} \cite{chen2018neuralode}, in which the hidden state
$x(t) \in \mathbb{R}^{D}$ is not updated by a discrete recurrence but defined as
the solution of
\begin{equation}
  \frac{\mathrm{d}x(t)}{\mathrm{d}t} = f\!\left(x(t), I(t), t, \theta\right),
  \label{eq:node}
\end{equation}
where $I(t)$ is the input and $f$ a neural network with parameters $\theta$. The
state at an arbitrary instant is obtained by numerically integrating
Equation~\eqref{eq:node}, and the network is trained by differentiating through
the solver. The formulation makes the model naturally defined on a continuous
time axis, but offers no guarantee on the asymptotic behaviour of the state.

A more stable formulation is that of \emph{continuous-time recurrent networks}
\cite{funahashi1993ctrnn}, which add to the learnt term an exponential decay
towards equilibrium,
\begin{equation}
  \frac{\mathrm{d}x(t)}{\mathrm{d}t}
    = -\frac{x(t)}{\tau} + f\!\left(x(t), I(t), t, \theta\right),
  \label{eq:ctrnn}
\end{equation}
in which $\tau$ is a time constant governing the speed at which the autonomous
system reaches equilibrium. The time constant is, however, a fixed parameter:
every unit forgets at the same speed, whatever it is observing.

The contribution of liquid cells consists in making that speed \emph{dependent
on the input and on the state}. The three cells compared in this work realise
the idea in three different ways: LTC by explicitly modulating the time constant
of the differential equation, CfC by replacing numerical integration with a
closed-form approximation of its solution, and LRCU by making variable the
capacitance of the equivalent circuit from which the equation derives. A
systematic comparison of the family against classical recurrent networks, on
accuracy, memory footprint and generalisation, is reported in
\cite{zong2025lnnsurvey}; the same line of research has since been extended to
structured state-space models \cite{hasani2023liquids4}, which are not employed
in this work.

\subsection{Liquid Time-Constant (LTC)}
\label{sec:cell-ltc}

The LTC cell \cite{hasani2021ltc} originates in an observation about the
dynamics of non-spiking neurons: the membrane potential evolves as a first-order
linear system, while the synaptic currents entering it are, in the steady state,
a sigmoidal nonlinearity of the state of the presynaptic neurons, multiplied by
the difference between the current potential and a reference potential $A$.
Substituting that term into Equation~\eqref{eq:ctrnn} yields
\begin{equation}
  \frac{\mathrm{d}x(t)}{\mathrm{d}t}
    = -\left[\frac{1}{\tau} + f\!\left(x(t), I(t), t, \theta\right)\right] x(t)
      + f\!\left(x(t), I(t), t, \theta\right) A .
  \label{eq:ltc}
\end{equation}

The difference with respect to Equation~\eqref{eq:ctrnn} is substantial. The
network $f$ no longer determines only the derivative of the state, but appears
also in the coefficient multiplying the state itself: it therefore acts as a
variable time constant of the system,
\begin{equation}
  \tau_{\text{sys}}
    = \frac{\tau}{1 + \tau f\!\left(x(t), I(t), t, \theta\right)} ,
  \label{eq:tausys}
\end{equation}
and it is from this dependence that the adjective \emph{liquid} derives. Each
unit of the hidden state may thus specialise on a time scale of its own and
modify it according to what it observes --- reacting quickly to an input that
changes, and retaining for a long time an information that does not --- whereas
Equation~\eqref{eq:ctrnn} would impose the same speed in both cases.

The model enjoys two properties proved in \cite{hasani2021ltc} that motivate its
use in a control loop. The first is that the time constant remains bounded,
$\tau/(1 + \tau W) \le \tau_{\text{sys}} \le \tau$, where $W$ collects the
weights entering the unit under consideration; the second is that the state of
each unit remains confined between $\min(0, A_{\min})$ and $\max(0, A_{\max})$
over any finite interval, whatever the magnitude of the input. The second
property is of greater practical relevance: the output of the recurrent core
cannot diverge even when the observation contains anomalous values, a condition
that in a real system occurs whenever a source of odometry degrades.

Equation~\eqref{eq:ltc} admits no analytical solution and constitutes a
\emph{stiff} system of equations, for which a Runge--Kutta integrator would
require a prohibitive number of steps. The authors therefore propose a
\emph{fused} solver, which replaces with the updated value only the occurrences
of the state that appear linearly, and can thus be solved symbolically,
obtaining the stability of the implicit Euler method at the cost of the explicit
one:
\begin{equation}
  x(t + \Delta t)
    = \frac{x(t) + \Delta t\, f\!\left(x(t), I(t), t, \theta\right) \odot A}
           {1 + \Delta t\left[1/\tau + f\!\left(x(t), I(t), t, \theta\right)\right]} .
  \label{eq:ltc-fused}
\end{equation}

\paragraph{In the implementation.}
Equation~\eqref{eq:ltc-fused} is realised as written, with
$f = \sigma(W_{\text{in}} u + W_{\text{rec}} x + b)$, with $\tau$ obtained as
$\operatorname{softplus}(\tau_{\text{raw}})$ so as to guarantee its positivity,
and with $A$ learnt. The step $\Delta t$ is fixed when the cell is constructed
and repeated a configurable number of times per input step, following the
\emph{unfolding} scheme prescribed by the authors. The initialisation of
$\tau_{\text{raw}}$ places the time constants on an interval comparable with the
rate of the data, so that the cell integrates meaningfully from the first step.

\subsection{Closed-form Continuous-time (CfC)}
\label{sec:cell-cfc}

The use of a numerical solver, however efficient, remains the principal cost of
continuous-time cells: every forward step requires one or more evaluations of
the differential equation, and the cost is multiplied by the length of the
sequence. The CfC cell \cite{hasani2022cfc} removes that cost by replacing the
integration with a \emph{closed-form approximation} of the solution of
Equation~\eqref{eq:ltc}, for which no analytical expression was previously
known. The authors report a speed-up of between one and five orders of
magnitude, in training and in inference, with respect to the equivalent
solver-based models.

The approximate solution takes the form
\begin{equation}
  x(t) = \sigma\!\left(-f(x, I; \theta_f)\, t\right) \odot g(x, I; \theta_g)
       + \left[1 - \sigma\!\left(-f(x, I; \theta_f)\, t\right)\right] \odot h(x, I; \theta_h),
  \label{eq:cfc}
\end{equation}
where $f$, $g$ and $h$ are three neural networks and $\sigma$ is the sigmoid.
The reading of Equation~\eqref{eq:cfc} is immediate: $g$ and $h$ describe
respectively the initial and the asymptotic behaviour of the unit, while the
sigmoid interpolates between the two as a function of the elapsed time, with a
transition speed determined by $f$, that is by the input and the state. Time
appears here \emph{explicitly} as an argument, and not as a discretisation step:
it is this property that makes the cell suited to a stream of observations at an
irregular rate, since a long interval and a short one produce by construction
two different values of the same gate.

\paragraph{In the implementation.}
The three networks share a common backbone, following the realisation proposed
by the authors: a single layer $s = \tanh(W [\,u, x\,] + b)$ produces a hidden
vector that three linear heads read, so that the cost of the cell remains close
to that of a traditional recurrence. The gate takes the form
$\gamma = \sigma\!\left(-\left(w_\tau + f(s)\right)\Delta t\right)$, in which
$w_\tau$ is a per-unit learnt time constant, and the updated state is
$x' = \gamma \odot g(s) + (1 - \gamma) \odot h(s)$.

The initialisation of $w_\tau$ deserves a remark, since it determines whether
the cell behaves as a continuous-time model at all. Set to zero, with a rate of
the order of a tenth of a second, the argument of the gate starts within a
hundredth of zero: the sigmoid pins the gate at $0.5$, the cell degenerates into
a constant blend of two heads, and the time constant for which it exists plays
no part. The initialisation adopted is therefore log-uniform over an interval of
time constants physically sensible for the problem --- $w_\tau$ is a rate, hence
the reciprocal of a time --- and places the gate in a usable region from the
first step.

\subsection{Liquid Resistance-Capacitance Unit (LRCU)}
\label{sec:cell-lrcu}

The third cell descends from the same equivalent electrical circuit from which
LTC moves, but revises two of its simplifications. The first concerns
saturation: in Equation~\eqref{eq:ltc} the synaptic conductance is not bounded
from above, whereas in a real circuit it is; saturated-conductance models
reintroduce the bound. The second, which constitutes the contribution proper to
\emph{liquid-resistance liquid-capacitance} networks \cite{farsang2024lrc},
concerns the membrane capacitance: in LTC it is constant, and only the
resistance varies with the state. Making the capacitance variable as well yields
\begin{equation}
  \frac{\mathrm{d}x}{\mathrm{d}t}
    = \varepsilon(w) \odot \left[-\,\sigma(f)\odot x + \tau(u) \odot e_{\ell}\right],
  \label{eq:lrc}
\end{equation}
where $\sigma(f)$ is the forget conductance, $\tau(u)$ the update conductance
towards the resting potential $e_{\ell}$, and $\varepsilon(w)$ is the
\emph{liquid elastance}, the reciprocal of the capacitance, itself a function of
the input and of the state. The third factor acts on the whole derivative: it
does not establish \emph{towards what} the unit evolves, but \emph{how quickly}
it does so, and is therefore a second mechanism of temporal modulation added to
the one already present in LTC.

The authors observe that integrating Equation~\eqref{eq:lrc} with an explicit
Euler method of order one and a single unfolding step yields a discrete
recurrence --- the LRCU unit --- structurally akin to a gated RNN, in which
$\varepsilon$ plays the role of the update gate of a GRU \cite{cho2014gru},
with in addition the conductance structure of the circuit from which it derives.
It is in this form that the cell is employed in the present work, since it
thereby retains the cost of a traditional recurrence.

\paragraph{In the implementation.}
The update is
\begin{equation}
  x' = x + \Delta t\, \varepsilon \odot
       \left(-f \odot x + z \odot e_{\ell}\right),
  \label{eq:lrcu}
\end{equation}
with the three quantities computed as
\begin{align}
  f           &= \operatorname{softplus}(g_{\max}) \odot \sigma\!\left(W^{\text{in}}_f u + W^{\text{rec}}_f x + b_f\right), \\
  z           &= \operatorname{softplus}(k_{\max}) \odot \sigma\!\left(W^{\text{in}}_z u + W^{\text{rec}}_z x + b_z\right), \\
  \varepsilon &= \sigma\!\left(W^{\text{in}}_\varepsilon u + w_\varepsilon \odot x + v_\varepsilon\right).
\end{align}
The forget and update conductances are therefore sigmoids scaled by a learnt
maximum conductance, according to the saturated form discussed above, while the
elastance has a diagonal rather than a full recurrence. The cell holds three
input matrices and two complete recurrent matrices, and is for this reason the
most expensive of the three at equal state width: the comparison of
Chapter~\ref{ch:results} is to be read bearing that asymmetry in mind, as
reported in Table~\ref{tab:architectures}.

\subsection{Time, and the asymmetry it introduces}
\label{sec:cell-time}

The three cells treat time in two distinct ways, and the difference must be
declared because it bears on the comparison. In the CfC the elapsed time appears
explicitly in Equation~\eqref{eq:cfc} as the argument of the gate: the cell
therefore receives the \emph{real} $\Delta t$ between two consecutive
observations and adapts its transition to it. In LTC and LRCU time appears
instead as the discretisation step of a numerical integration
(Equations~\eqref{eq:ltc-fused} and~\eqref{eq:lrcu}), fixed when the cell is
constructed; the real $\Delta t$ reaches them as an additional component of the
input, which is the reason why their observation vector is \num{42} values wide
rather than \num{41}.

The asymmetry is a consequence of the formulation and not a choice of the
implementation. Since the rate of the observations is not constant
(Section~\ref{sec:observation-recording}), the constant step is inexact for a
non-negligible fraction of the samples, and the same quantity reaches the CfC in
a more informative form than that in which it reaches the other two. Supplying
LTC and LRCU with a per-sample integration step would alter their update rule
--- and, in the case of LTC, would invalidate the stability bounds recalled in
Section~\ref{sec:cell-ltc}, which presuppose a fixed step --- so that the
comparison would be between a cell and a variant of it, rather than between the
cells proposed in the literature.

\section{Wiring: Dense and NCP}
\label{sec:lnn-wiring}

The second choice concerns the connectivity of the recurrent core. In the dense
variants every weight matrix is full and the cascade is direct: fusion, first
cell, second cell, head.

The alternative variants draw on \emph{Neural Circuit Policies}
\cite{lechner2020ncp}, in which a continuous-time cell is embedded in a sparse
wiring derived from the nervous system of the nematode \emph{Caenorhabditis
elegans}. The original wiring organises the neurons in four groups ---
sensory neurons, interneurons, command neurons and motor neurons --- with sparse
connections between adjacent groups, recurrence concentrated in the command
group, and synapses whose sign, excitatory or inhibitory, is established by the
structure of the circuit rather than learnt. The authors show that an end-to-end
driving policy can be realised with \num{19} control neurons and \num{253}
synapses, and attribute to such compactness three advantages: interpretability
of the learnt policy, robustness with respect to perturbations of the input, and
a reduction of the parameter space.

Of that approach the present work adopts the two elements that can be applied to
a recurrent core of fixed width without redesigning its topology. The first is
\emph{signed sparsity}: the matrices of each cell are multiplied by a mask
generated at initialisation, which removes a fraction of the connections and
assigns to each of the remaining ones a constant sign that training cannot
reverse. The density is \num{0.6} on the input connections and \num{0.8} on the
recurrent ones. The masks are registered as buffers and are therefore part of
the checkpoint, so that a reloaded network resumes exactly the connectivity with
which it was trained. The second is the presence of \emph{paths that bypass one
level}, analogous to the diagonal connections of the biological circuit: from
the input to the second cell, and from the first cell to the head, both at
density \num{0.4}. The former is supported natively by the cells; the latter,
since the head is a linear layer and not a cell, is realised as a separate
masked additive path, leaving the main weights of the head dense.

The four-group topology of the original wiring is not adopted. The reason is
that the architecture described in this chapter places the recurrent core
between a convolutional encoder and a calibrated head, both dense and common to
all variants: introducing a sensory group and a motor group would mean modifying
the input and output paths as well, and the comparison would no longer isolate
connectivity alone. What the sparse variants of this work measure is therefore
the contribution of signed sparsity and of the bypass paths, not that of the
complete circuit proposed in \cite{lechner2020ncp}; the designation NCP is
retained for continuity with the literature and with the implementation, but is
to be understood in this restricted sense.

The dense variants have neither masks nor bypass paths, and are identical in
everything else --- encoder included --- to their sparse counterpart. The pair
therefore isolates exactly what sparsity and the bypass paths are worth.

\section{Sparse Convolutional Encoder}
\label{sec:sparse-encoder}

The encoder reduces the local voxel map to a descriptor of fixed length. It is a
four-level pyramid of \num{16}, \num{24}, \num{32} and \num{48} channels, each
preceded from the second onwards by a downsampling stage of stride~\num{2}. The
convolutions are \emph{submanifold} \cite{graham2018submanifold}, that is they
keep the computation anchored to the occupied voxels alone instead of densifying
the grid: it is the property that makes tractable an input in which almost all
cells are empty. Every level
is pooled per sample by a reduction over the active sites, and the four
descriptors are concatenated into a vector of \num{200} channels.

Two normalisation choices deserve to be motivated, since neither follows from
the pyramid itself. The first is the use of \emph{group normalisation} in place
of \emph{batch normalisation}. Batch statistics would here be computed over the
active voxels of whatever samples share a batch, and the number of active voxels
varies appreciably from one recording to another: the statistics would therefore
depend on the composition of the batch rather than on the scene, and the
descriptor produced for a fixed input would move from one epoch to the next.
Group normalisation removes that dependence, since it normalises within each
sample. The second is a \emph{layer normalisation} applied to the pooled output.
The dense branch reaches the fusion already standardised to unit variance,
whereas the pooled descriptor carries the scale of the pooling: without a
normalisation the point-cloud branch would dominate the fusion by scale alone,
that is for a reason unrelated to its information content.

The encoder accepts a different spatial shape at every call, because the
convolutions operate on the offsets between active voxels and not on absolute
indices. It is not, however, invariant to the voxel pitch: a kernel of fixed
size covers a different physical extent at a different resolution, which is the
reason why the split of the corpus stratifies on that parameter
(Section~\ref{sec:sl-split}).

\section{Observation Normaliser and Model Embedding}
\label{sec:aux-modules}

\paragraph{The observation normaliser.}
The \num{41} dense values enter the network in a fixed order --- goal, odometry,
state --- and are transformed by a fixed affine map whose mean and standard
deviation are \emph{estimated from the data before training} rather than learnt.
For the LTC and LRCU variants the values are \num{42}, because the real elapsed
time is concatenated by the trainer (Section~\ref{sec:lnn-taxonomy}).

The order of the fields is part of the contract and not an internal convention:
getting it wrong raises no error, because the widths continue to match, and
simply produces wrong predictions.

\paragraph{The robot-model embedding.}
The model identifier is a categorical variable written as the last value of the
state block. Left there, it reaches the network as a large continuous number,
which the normaliser reduces to a couple of arbitrary values and which says
nothing about the identifiers being unordered categories. The embedding assigns
to each robot model a vector of its own, of eight dimensions. The table of
identifiers is a buffer, so every checkpoint carries its own mapping, and one
row is reserved for a model never seen in training.

\section{Output Calibration}
\label{sec:output-calibration}

The final stage of the network --- a linear head followed by a recalibration ---
produces the six components of the command directly in physical units. Nothing
saturates them against the kinematic limits of the platform: what the network
says is what goes out on the command channel.

The recalibration exists for a precise reason. A head trained with a squared
error learns the conditional mean of the action; where the demonstration is
multimodal for the same observation --- an alternation between standstill and
travel near the velocity limit, with little in between, is enough --- that mean
falls between the modes and belongs to neither. The output of such a head is
therefore systematically narrower than the commands it was trained on. In open
loop this is an error like any other; in closed loop, where the command becomes
the next observation, it is a robot that travels at a fraction of the
demonstrated speed and takes correspondingly longer to reach its goal, so that
the shortfall compounds along the trajectory. The amplitude metrics of
Section~\ref{sec:metrics-amplitude} exist to measure it.

The calibration is the least-squares line from prediction to command, per
component,
\begin{equation}
  a_i = \frac{\operatorname{cov}(\hat{y}_i, y_i)}{\operatorname{var}(\hat{y}_i)},
  \qquad
  b_i = \mathbb{E}[y_i] - a_i\,\mathbb{E}[\hat{y}_i],
\end{equation}
estimated on the validation set at every epoch and stored as buffers. Being the
best affine map of the model's own output, it can never do worse than the best
constant predictor. The gain can be rewritten as
$\rho_i\,\sigma_{y_i}/\sigma_{\hat{y}_i}$, and therefore restores amplitude only
in proportion to how much the model actually tracks the command. The alternative
of enforcing a match of the \emph{spreads} --- a gain equal to
$\sigma_{y_i}/\sigma_{\hat{y}_i}$, which makes the output exactly as wide as the
demonstration --- ignores the correlation and amplifies mostly noise wherever
the model tracks weakly, and may therefore leave a weakly tracked component
worse than the constant predictor, which the least-squares line cannot do.

Two properties must be stated. The first is that the model must be
\emph{trained} with this stage at the identity: it lives in the forward pass, so
a gain left in place during the following epoch is optimised through, and the
network learns to shrink its raw output by exactly the factor the gain restores.
The second is what the calibration cannot do: being monotone per component, it
cannot separate two modes that the model reports as a single number. It makes
the answer the right size, not more correct.

\section{Auxiliary Geometry Head}
\label{sec:geometry-head}

Imitation alone teaches the encoder very little: the demonstrator flies a path
planned on the true geometry of the world, so its command depends on the local
map only indirectly, and the map matters only in the instants in which the
vehicle deviates for an obstacle. The geometry surrounding the robot, however,
is present in every frame and requires no labels: the cloud itself states where
the obstacles are.

The auxiliary head reads the descriptor of the encoder and predicts that
geometry: the clearance in eight azimuth sectors and the height above the
ground. The targets are measured on the cloud of each frame with the same
functions the perception bridge uses on board, so that the encoder is taught
exactly the quantity the reinforcement learning reward prices
(Section~\ref{sec:rl-reward}). The distances are divided by the map radius of
the episode before the loss, so that grids of different extent share a single
scale and a sector with nothing in range reads exactly~\num{1}; entries the
bridge could not measure are left out of the loss rather than counted as zero.

Two design choices keep the head in its role. It is deliberately narrow --- a
single hidden layer --- because its purpose is to shape the encoder, and a head
wide enough to decode a poor descriptor on its own would absorb the gradient
instead of transmitting it. And it is an auxiliary task and not part of the
policy: it lives beside the network in the trainer and is saved in a file of its
own, so that the checkpoint inference loads keeps exactly the structure of the
network it executes.

\section{PyTorch Module Implementation}
\label{sec:module-implementation}

The formulations of the preceding sections are realised in PyTorch on two
levels. The first collects the \emph{cells} --- the three recurrent units, in
their dense and sparsely wired versions --- and the stateless modules the
architecture employs: encoder, normaliser, embedding, calibration and auxiliary
head. The second collects the \emph{networks}, each of which composes those
modules into a complete architecture. The separation mirrors that of the
chapter: the first level contains what is drawn from the literature, the second
what was designed for this problem.

\subsection{The interface of the cells}

All cells derive from a base class that establishes two things: the signature of
the forward step, and the handling of the recurrent state. What the base class
deliberately does \emph{not} contain is the update equation, not even in those
parts in which two cells would coincide: the equation is the element that
defines the cell, and remains readable in the file of the cell that realises it.
It follows that replacing LTC with CfC inside a network requires no change to
the network.

The state lives in the cell as an ordinary tensor and not as a registered
buffer. The distinction has a practical consequence: the state does not enter
the dictionary of the weights, so that a checkpoint contains the learnt
parameters and not the memory of one particular sequence. It is allocated at the
first call, with the size of the incoming batch, and is manipulated through four
operations that training employs at distinct moments:

\begin{itemize}
  \item \texttt{get\_state} returns a copy \emph{detached} from the
        differentiation graph. It is the initial condition from which the
        following window resumes in truncated backpropagation through time
        (Section~\ref{sec:sl-sequences}): without the detachment, the current
        window would carry with it the graph of the preceding one.
  \item \texttt{set\_state} restores a copy so obtained, or clears the state
        when it receives a null value.
  \item \texttt{detach\_state} severs the graph after every optimiser step; in
        its absence the graph would grow without bound from one mini-batch to
        the next.
  \item \texttt{reset\_state} zeroes the state, either for the whole batch or
        only for the positions indicated by a boolean mask.
\end{itemize}

The last operation is the one that makes the cells usable in the vectorised
environment of Chapter~\ref{ch:training}, in which the episodes of the robots of
the fleet end at different instants: only the positions corresponding to
finished episodes must restart from a null state, while the others continue with
their own memory.

\begin{lstlisting}[language=Python, caption={Selective reset of the recurrent state.}, label={lst:reset}]
def reset_state(self, mask=None) -> None:
    if self._x_t is None:
        return
    if mask is None:            # whole batch
        self._x_t = None
    else:                       # only the flagged entries
        # clone first: the tensor may still belong to the
        # autograd graph, and must not be written in place
        self._x_t = self._x_t.clone()
        self._x_t[mask] = 0.0
\end{lstlisting}

\subsection{The sparse wiring}

The wiring described in Section~\ref{sec:lnn-wiring} is realised as a mask
generated once when the cell is constructed and registered as a buffer, so that
it is saved together with the weights and a reloaded network resumes the
connectivity with which it was trained. The mask takes values in
$\{-1, 0, +1\}$: the null value removes the connection, the sign makes it
permanently excitatory or inhibitory. It is applied by multiplying it into the
matrix it governs, so that the gradient reaches only the connections that are
present and the sign cannot be reversed by training.

\begin{lstlisting}[language=Python, caption={Generation of the signed sparse mask.}, label={lst:mask}]
def signed_mask(shape, density: float) -> torch.Tensor:
    active = torch.bernoulli(torch.full(shape, density))
    sign   = torch.where(torch.rand(shape) < 0.5, -1.0, 1.0)
    return active * sign
\end{lstlisting}

\subsection{The composition of a network}

Every network builds in its own constructor everything that defines it ---
encoder, cells, heads --- rather than inheriting it from a hierarchy. The choice
has a cost in repetition and a precise benefit: the variant that runs on board
the robot is self-contained, and the inference node of Chapter~\ref{ch:pipeline}
carries its file and the cells it employs, without depending on the rest of the
training repository.

From the common base class the networks inherit instead two things. The first is
the propagation of the state: the four operations described above are applied in
sequence to all the cells of the network, so that whoever trains manipulates the
state of the network and not that of the individual units. The second is a set
of properties by which the network \emph{declares itself} --- whether it is
recurrent, whether it requires the elapsed time as an argument, whether it
consumes a point cloud, whether it uses the robot-model embedding. Both trainers
query those properties in order to build the batch and to compose the call, and
neither contains a check on the type of the network: adding an architecture
therefore does not entail modifying the trainers.

The forward step reflects the architectural scheme directly. The three input
branches are gathered into a single fusion vector, which traverses the two cells
in cascade; in the sparse variants it also reaches the second cell as a bypass,
while the output of the first is added to the pre-activation of the head through
its own masked path; the result is finally emitted through the calibration of
Section~\ref{sec:output-calibration}.

\begin{lstlisting}[language=Python, caption={Forward step of a sparse variant, from the three sources to the command.}, label={lst:forward}]
# the three branches, fused into one vector
x  = self.fuse(obs, coords, feats, spatial_shape, model_id, encoded)

h1 = self.cfc_1(x, dt)
h2 = self.cfc_2(h1, dt, skip=x)    # input -> cell 2, past cell 1

out = self.hidden_out(self.tanh_1(h2)) \
    + h1 @ (self._skip_hidden_W * self._skip_hidden_mask)

return self.emit(self.tanh_2(out))   # head -> calibration
\end{lstlisting}

\subsection{The selection of the architecture}

Every architecture --- the six candidates, the stateless control and the vector
twins of Section~\ref{sec:networks-implementation} --- is declared in a single
table, which associates with each of them the class that realises it and the
hyper-parameters that define it. A
construction function consults the table and returns the requested network; both
trainers and the measurement tools go through that function. The tabular form,
in place of a chain of conditional branches, has two intended consequences: the
dimensions of all the architectures are readable side by side, which makes
verifiable the claim that the comparison holds everything equal except the cell
and the wiring; and the addition of an architecture is a single line, which
cannot forget a branch.

Declared in the same place is the set of architectures that integrate with a
fixed step, that is the four LTC and LRCU variants
(Section~\ref{sec:cell-time}). Membership of that set determines two things that
must necessarily agree: the width of the input layer of the network, increased
by one component, and the decision of the trainer to concatenate the elapsed
time to the observation before the call. Since both read the same declaration,
they cannot diverge --- and a divergence would be silent, given that a network
built for \num{41} inputs and fed with \num{42} values would raise a shape error
only if the two widths did not happen to coincide.

\begin{lstlisting}[language=Python, caption={The architecture table and the declaration of the fixed step.}, label={lst:table}]
# the cells that integrate with a fixed step, declared once
DT_FOLDED_INTO_OBS = frozenset({
    Model_types.SPARSECONV_LTC_NCP,
    Model_types.SPARSECONV_LRCU_NCP,
    Model_types.SPARSECONV_LTC,
    Model_types.SPARSECONV_LRCU,
})

ARCHITECTURES = {
    Model_types.SPARSECONV_CFC_NCP: (SparseConv_CfC_NCP_Network,
        dict(connections=0.6, fusion_dim=192, head_dim=128,
             cfc1_out_dim=80, cfc2_out_dim=80,
             cfc1_hidden_dim=80, cfc2_hidden_dim=80)),
    Model_types.SPARSECONV_LTC_NCP: (SparseConv_LTC_NCP_Network,
        dict(connections=0.6, fusion_dim=192, head_dim=128,
             ltc1_out_dim=80, ltc2_out_dim=80)),
    # ... one entry per architecture
}

# the extra input goes through the same declaration
in_features = obs_dim + 1 if arch in DT_FOLDED_INTO_OBS \
              else obs_dim
\end{lstlisting}

\subsection{The checkpoint}

At the end of every epoch training writes a directory containing the weights of
the network, those of the auxiliary head in a separate file, the state of the
optimiser and of the scheduler, and a metadata file. The separation of the
auxiliary head is not incidental: it does not belong to the policy
(Section~\ref{sec:geometry-head}), and keeping it out of the weight file allows
the inference node to load the latter strictly, that is rejecting any parameter
that does not belong to the network it intends to execute.

The metadata file declares the architecture, the options with which its encoder
was built and the observation contract the network assumes, besides the
parameters of the training session and the position from which it may resume.
The principle is that \emph{nothing infers an architecture from the shape of its
weights}: whoever reloads a checkpoint --- the reinforcement training that
starts from the imitated weights, the measurement tools, the on-board node ---
rebuilds the network from the declaration and loads the weights into it, rather
than attempting to guess its structure. An inference from the shapes would work
in most cases and fail silently in the remaining ones, since two variants
differing only in the wiring have matrices of the same dimensions.

\section{Network Architectures Under Comparison}
\label{sec:networks-implementation}

The comparison of Chapter~\ref{ch:results} concerns the six architectures that
consume the point cloud, obtained by crossing the three cells with the two
wirings, together with a seventh network that is not a candidate but a control
(Table~\ref{tab:architectures}).

\begin{table}[H]
  \centering
  \caption{The six architectures under comparison and the stateless control:
           what distinguishes them is the cell and the wiring; the parameter
           count is a consequence, not a choice. For the control the last column
           is the share of the two linear layers that take the place of the
           cells.}
  \label{tab:architectures}
  \begin{tabular}{@{}lccrr@{}}
    \toprule
    \textbf{Network} & \textbf{Input} & \textbf{Real $\Delta t$} & \textbf{Parameters} & \textbf{in the cells} \\
    \midrule
    \met{SparseConv_CfC_NCP}  & 41 & yes & 282\,430 & 31.6\,\% \\
    \met{SparseConv_LTC_NCP}  & 42 & no  & 243\,902 & 20.7\,\% \\
    \met{SparseConv_LRCU_NCP} & 42 & no  & 331\,582 & 41.6\,\% \\
    \met{SparseConv_CfC}      & 41 & yes & 256\,830 & 28.7\,\% \\
    \met{SparseConv_LTC}      & 42 & no  & 218\,302 & 16.1\,\% \\
    \met{SparseConv_LRCU}     & 42 & no  & 275\,262 & 33.4\,\% \\
    \midrule
    \met{SparseConv_MLP} (control) & 42 & no & 205\,182 & 10.7\,\% \\
    \bottomrule
  \end{tabular}
\end{table}

\paragraph{The control.}
The seventh row is the six with the recurrence taken out: the same encoder,
fusion, stage widths, head and output stage, with two plain linear layers where
the cells were, and therefore no state carried from one step to the next. It
answers the question the six cannot answer among themselves --- what the memory
is worth on this data --- and it sets the bar a recurrent variant has to clear:
one that does not beat the control has not earned its state. Since it has no
integration step, the elapsed interval reaches it as it reaches LTC and LRCU, as
a trailing component of the observation.

\paragraph{The vector twins.}
Each of the seven networks has a twin that reads the dense observation alone:
the same cells and wiring, the same bypass paths, the same head and output
stage, with the encoder removed. The twins are what the synthetic validation
tasks of Section~\ref{sec:exp-protocol} are run on --- the sine wave for the
supervised stage, LunarLander for the reinforcement stage --- neither of which
has a point cloud, so that what a cell does there is evidence about the same
cell inside the network that flies. Three differences are deliberate. There is
no encoder, and the fusion takes the observation alone. The widths are those the
toy tasks were sized with, shared by all seven twins: an input layer of
\num{128}, two recurrent stages of \num{64} and a head of \num{128}, against
the $192/80/80/128$ of the spatial networks. And the input layer, the head and
the output projection are initialised orthogonally, the last with a small gain,
because the twins are trained from scratch, where a near-zero output is the
customary starting point of a policy; the spatial networks keep the default
initialisation because their weights come from imitation. The elapsed interval
reaches each twin exactly as it reaches its spatial counterpart.

\paragraph{What is held equal.}
The observation and its normalisation, the encoder and its configuration, the
width of the fusion, the width of the two recurrent states, the head, the output
stage with its calibration, and every training setting.

\paragraph{What is not, and is to be read alongside the result.}
The parameter count and the way in which time reaches the cell. The spread of
the parameter counts is intrinsic: at equal state width the three cells have
very different weight densities, and LRCU carries nearly three times as many as
LTC. Forcing equal counts would mean unequal widths, which changes what is being
compared; comparing at equal width and reporting the counts is the more honest
of the two alternatives.

\paragraph{Sizing.}
The widths were chosen for the problem rather than for the corpus at hand, and
are held equal across the networks compared so that the comparison does not
become a comparison of capacities. Whether the capacity so chosen is the right
one is a question for the experiments rather than for this chapter: the relation
between validation and training loss (Section~\ref{sec:metrics-loss}) and the
occupancy of the recurrent state (Section~\ref{sec:metrics-health}) are logged
for that purpose, and Chapter~\ref{ch:results} reads them. What belongs here is
the direction of the trade-off: a narrower recurrent state would remove the
headroom that reinforcement learning is expected to use for exactly the
behaviour the supervised phase cannot teach.

\paragraph{What to judge on.}
Not on the loss alone. The metrics defined in
Section~\ref{sec:metrics-supervised} separate what the loss conflates: whether
the model beats echoing the velocity it has already been given, whether it
commands as much as the demonstration did --- which a squared error rewards
\emph{not} doing --- and how it behaves separately on the two classes of robot,
which differ enough that an average hides one of them.
